% 01-problem.tex: задача, которую пытается решить Polaris.
\section{Задача}
\label{sec:problem}

Человек в Соединённых Штатах носит с собой от шести до восьми удостоверяющих документов, которые
друг с другом не разговаривают: водительское удостоверение, паспорт, карточку социального
страхования, Real ID, регистрацию избирателя, страховую карту. Каждый из них есть отдельный
предмет, подписанный отдельным органом по отдельному стандарту, без общего пути отзыва и без
общего следа для аудита. Версия 1 настоящего отчёта бралась смоделировать их объединение в одну
действующую запись удостоверения на человека, проверяемую через события с областью действия,
заданной контекстом. Это лёгкая половина задачи, и она по-прежнему задаёт её очертания.

Тяжёлая половина состоит в том, что происходит, когда система работает. Идентификационная
система, добившаяся успеха, становится дверью ко всему, а дверь есть место, где
сосредоточивается власть. Отсюда следуют три режима отказа, и каждый слой Polaris существует,
чтобы ответить на один из них.

\begin{description}
  \item[Система лжёт о прошлом.] Оператор или инсайдер с доступом к базе данных отзывает токен, а
    позже стирает запись; выдаёт удостоверение по процедурам, о которых обществу не говорили; или
    показывает одну историю аудитору и другую суду. Национальной идентификационной системе нельзя
    позволять задним числом переписывать то, что она сделала.
  \item[Система наблюдает.] Каждая проверка есть факт о том, где человек был и чем занимался.
    Журнал таких фактов есть хребет слежки, независимо от того, задумывал ли его кто-нибудь, и он
    остаётся им после того, как намерения оператора переменятся.
  \item[Система принуждает.] Держателя можно заставить: пограничник, работодатель или вор могут
    вынудить предъявить удостоверение, а дверь, которой требует каждая служба, способна вынудить
    население зарегистрироваться. Идентификационная система есть также и самый действенный
    инструмент исключения из всех когда-либо построенных, если список тех, кого в ней нет, можно
    получить по требованию.
\end{description}

Общепринятый ответ на все три есть политика: правила о сроках хранения, доступе, аудите и
согласии, применяемые теми людьми, которые систему и эксплуатируют. Конституция Polaris
(\code{MISSION.md}) держится позиции, что политика ничего не стоит как единственная опора
обещания, потому что люди, эксплуатирующие систему, суть ровно те, кого обещание и призвано
связать. Обещание должно быть закреплено там, где его нельзя обойти разговором: в схеме, в
триггерах, в уникальных индексах, в ограничениях целостности, в подписях, которые любой может
проверить, ни о чём систему не спрашивая, и в журналах, за которыми наблюдают посторонние.
Приложение не есть граница безопасности. Ею является база данных, а слои вокруг базы данных
существуют затем, чтобы даже её собственные утверждения могли проверить посторонние.

\subsection{Почему постквантовость с первого дня и что этой фразе позволено означать}

Версия 1 доказывала через неравенство Моски, что горизонт секретности идентификационных данных
измеряется поколениями, что миграция национальной системы занимает десятилетие и что система
удостоверений, развёрнутая на классической криптографии в 2026 году, есть поэтому отложенная
компрометация, а не работающая система. Этот довод не изменился и кратко изложен в
приложении~\ref{app:from-v1}. На \code{1.0.0-rc.7} алгоритмом подписи по умолчанию является
ML-DSA-65 (FIPS 204), рядом с ним принимается ML-DSA-87, система доказательств с нулевым
разглашением построена на хеш-функциях и не требует доверенной настройки, а публичная кромка TLS
и один из двух внутренних переходов согласуют гибридный постквантовый обмен ключами.

То, что репозиторий об этом утверждает, сузилось со времён версии 2, и само это сужение есть один
из предметов настоящей версии. Слово \emph{постквантовый} точно в отношении алгоритмов. Оно не
является утверждением, которое репозиторий вправе делать о безопасности, потому что стойкость
ML-DSA опирается на решёточные допущения, которые могут пасть от математики, а не от железа, и
потому что путь разработки по умолчанию подписывает помеченной заглушкой, а вовсе не решёточной
подписью. Защитимое утверждение звучит так: \emph{алгоритмическая гибкость при проверяемом пути
миграции}. Алгоритм есть поле в схеме и никогда не константа, население можно переподписать по
новому алгоритму путём, который выполняется при каждой отправке, и во что обошёлся бы взлом,
записано. Раздел~\ref{sec:signatures} приводит довод и цену; инвариантная проверка отказывает
любой внешней поверхности, говорящей больше.

\subsection{Последовательность, которой следует настоящий документ}

Каждое из следующих предложений называет слой, и каждое истинно в отношении слоя, стоящего перед
ним.

\begin{itemize}
  \item Подлинность удостоверения не есть его нынешние полномочия. Подпись остаётся настоящей
    навсегда; полномочия под ней могут быть отозваны сегодня после обеда.
  \item Нынешние полномочия не есть приватность. Верификатор может узнать, что удостоверение
    действительно, и всё же узнать куда больше, чем ему требовалось.
  \item Приватность не есть свобода от принуждения. Система может не знать о человеке почти
    ничего и всё же оставаться дверью, которой требует каждая служба.
  \item Федерация не есть институциональная независимость. Два экземпляра одного и того же кода,
    запущенные одним и тем же автором, доказывают, что протокол работает, а не что его
    эксплуатируют независимые учреждения.
  \item Журнал прозрачности не есть защита от расщеплённого обзора. Одному наблюдателю могут
    показать частное ответвление истории, и он об этом не узнает.
  \item Зелёный тест не есть доказательство того, что тест задал верный вопрос.
  \item Система, проверившая саму себя, не есть система, которую проверил кто-то ещё. Каждый
    инструмент в репозитории прошёл мимо дефекта, который первая же сторонняя реализация
    отвергла с первого взгляда.
\end{itemize}

Последние два предложения дались проекту дороже всего. Раздел~\ref{sec:senses} о шестом, а
раздел~\ref{sec:verifiers} несёт седьмое.
